% Introduction. Architecture follows CH7 (question, taxonomy, results in
% advance, organisation) and CH2 (threads), not the four-preface CH8 template.

\section{Introduction}
\label{var:sec:introduction}

Every model carried through to a result in the preceding chapters holds its
rates fixed. A replicative unit is born at rate $\lambda$, dies at rate $\mu$
and triggers
catastrophe at rate $\delta$, and those three numbers are the same at the end
of the process as at the beginning. That assumption is what makes the
generating functions close, and it buys a great deal: the roots $a$ and $b$,
the burst-size law, the flooding criterion and the removal budget. It is also
the first thing an ecologist or an epidemiologist will question, because in
almost every setting one might want to describe, a rate is not a constant ---
resources run out, competitors adapt, and a host mounts an immune response.

This chapter asks what happens when a rate is allowed to move. The models
collected here were built at different times and stand at different stages of
completion, but they sort by one structural feature: \emph{what moves the
rate}. Four mechanisms recur.

\begin{figure}[!htbp]
\centering
\includegraphics[width=0.98\textwidth]{fig01_taxonomy/fig01}
\caption{Four ways a rate can move. Remaining room that is spent and never
returned yields a proper limit law; a potential that is restocked from outside
yields a cycle; resistance that drags a critical point yields reversal; a
rupture at rate $\delta_1 X_t$ yields two phases that compose, because the
burst-size law has already integrated the random time out.}
\label{var:fig:taxonomy}
\end{figure}

Where a rate is driven by the process's own state, the behaviour turns on
whether the feedback exhausts a quantity or replenishes it. Speciation under
logistic growth consumes room that is never returned, the integrated rate
converges, and the species count has a proper limit law. An accumulating
potential, restocked at a constant inflow, instead yields a cycle. Resistance
in the Pimentel model does something else again: it acts not on a level but
on a \emph{critical point}, dragging the state of no preferred direction after
the leading population and so restoring criticality. In the last model the
rate jumps at a single instant rather than bending continuously: the
intracellular phase ends at the rupture of the host cell, which the process
itself brings on at rate $\delta_1 X_t$, and the burst-size law has already
integrated that time out. That is why the two phases compose. The ecological
models occupy the first half of the chapter; a reader who wants only the
\yp{} material will find it in
\cref{var:sec:publicgood,var:sec:twophase}.

Four results carry the finished mathematics, and each is stated here so that
the derivations may be read knowing where they go.

\begin{itemize}[leftmargin=1.6em,itemsep=0.35em,topsep=0.35em]
\item For a pure birth process from one founder with per-capita rate
$\beta(t)$, the species count $S_t$ is geometric on $\{1,2,\dots\}$ with
parameter $\ee^{-\Bint(t)}$, where
$\Bint(t)=\int_0^t\beta(u)\,\dd u$
(\cref{var:prop:yule-geom}). The rate law enters the distribution only
through its integral.
\item Under the logistic rate the integral converges, so the process has a
proper limit law and not merely a finite late-time mean: $S_\infty$ is
geometric with mean $(K/N_0)^{\sigma K/\varrho}$
(\cref{var:cor:logistic-limit}), depending on the parameters through exactly
those two dimensionless groups.
\item In the Pimentel competition model the state at which the chain is
equally likely to step up as down is available in closed form
(\cref{var:eq:icrit}), and it is unstable. Two trajectories started from the
same state diverge exactly when they are near it.
\item Composing an intracellular birth--death--catastrophe process with an
extracellular birth--death process gives the host-to-host mean offspring
$\ex{Z}$ in closed form (\cref{var:eq:EZ}), as an explicit function of the
phase-1 catastrophe rate.
\end{itemize}

The Yule process, its geometric marginal and the method of characteristics
are \ChM{}'s; the logistic speciation model itself, with its mean, is already
a worked example there. The birth--death--catastrophe process, the roots $a$
and $b$ and the burst-size law are \ChCore{}'s and \ChDist{}'s. None of that
is re-derived here. What this chapter adds is the distribution under a moving
speciation rate, the moving critical point of the Pimentel chain, the
two-phase extinction probability, and the remaining specifications, stated
precisely enough to be taken up where they stop.

\begin{table}[H]
\centering
\small
\begin{tabular}{@{}c>{\raggedright\arraybackslash}p{0.15\textwidth}>{\raggedright\arraybackslash}p{0.35\textwidth}>{\raggedright\arraybackslash}p{0.21\textwidth}@{}}
\toprule
\S & Model & What moves the rate & Status \\
\midrule
\ref{var:sec:speciation} & Logistic speciation
  & remaining capacity $K-N(t)$ & closed form \\
\ref{var:sec:pimentel} & Pimentel competition
  & evolving resistance, itself driven by the state
  & closed-form critical point; simulated \\
\ref{var:sec:pimentelplus} & Pimentel+
  & fitness driven by dominance, reduced by species count & specified \\
\ref{var:sec:rrrr} & Rise, run, ruin, rejuvenation
  & potential $V$; catastrophe by scarcity $X_t/V$ & simulated \\
\ref{var:sec:plague} & Birth--plague
  & loss channel driven by the coupled component $Y_t$ & simulated \\
\ref{var:sec:publicgood} & Public good
  & $\lambda(t)=\lambda_0+\alpha P(t)$, $P$ made by the population
  & $K^\ast$ computed \\
\ref{var:sec:twophase} & Two phases of infection
  & switch at the rupture $t_c$; decaying $\mu_2(t)$ & closed form \\
\bottomrule
\end{tabular}
\caption{The seven models, and the mechanism by which a rate moves in each.
The last column is the shape of the chapter:
\cref{var:sec:speciation,var:sec:twophase} are its two mathematical ends; the
rows between them halt at a closed-form fragment, a simulation or a
specification. The birth--plague model is carried in two variants, so the
count of distinct specifications is eight.}
\label{var:tab:models}
\end{table}

\subsection{Notation}
\label{var:sec:notation}

Symbols are reused across models but never within one section, and never
against a thesis-wide meaning. The table records the assignment.

\begin{table}[H]
\centering
\footnotesize
\setlength{\tabcolsep}{4pt}
\begin{tabular}{@{}ll>{\raggedright\arraybackslash}p{0.56\textwidth}@{}}
\toprule
\S & Symbol & Meaning \\
\midrule
all & $\lambda,\mu,\delta$ & birth, death and catastrophe rates, as in
  \ChCore{} \\
& $a,b$ & roots of the characteristic quadratic, $b<1<a$
  (\cref{bdc:lem:roots}) \\
& $\Kb$ & burst size of a birth--death--catastrophe process \\
\midrule
\ref{var:sec:speciation}
 & $S_t$ & species count, $S_0=1$ \\
 & $N(t),K,N_0$ & total abundance, carrying capacity, initial abundance \\
 & $\varrho$ & logistic growth rate \\
 & $\sigma$ & speciation coefficient; $\beta(t)=\sigma(K-N(t))$ \\
 & $\xi$ & the constant $K/N_0-1$ \\
 & $\beta(t)$ & speciation rate per species; $\beta$ is this and nothing else \\
 & $\Bint(t)$ & integrated rate $\int_0^t\beta(u)\,\dd u$; $\eta$ is a constant floor added to $\beta$ in two of the alternative laws \\
\midrule
\ref{var:sec:pimentel},
\ref{var:sec:pimentelplus}
 & $A_t,B_t$ & counts of the two competing populations, $A_t+B_t=N$ \\
 & $N$ & capacity of the enclosure, an integer \\
 & $R_a,R_b$ & resistance of each population \\
 & $p_a,p_b$ & blocking probabilities, at most one non-zero \\
 & $\varepsilon,\alpha$ & evolution rate (the step size of the resistance update) and the exponent sharpening that drive near a boundary \\
 & $i_c$ & the critical state \cref{var:eq:icrit} \\
 & $f_a,f_b$ & fitness, the Pimentel+ recasting of $R_a,R_b$ \\
 & $S^{(a)}_t,S^{(b)}_t$ & per-population species counts; $\zeta$ couples dominance to speciation and $\gamma$ penalises fitness per species \\
\midrule
\ref{var:sec:rrrr}
 & $X_t,V(t)$ & count of replicative agencies, and accumulated potential \\
 & $\phi,c,\omega$ & potential inflow, potential consumed per agency, revival rate out of the empty state \\
\midrule
\ref{var:sec:plague}
 & $X_t,Y_t$ & uninfected and infected counts \\
 & $\chi$ & contact-infection coefficient \\
 & $V(t),\phi$ & vitality and its inflow rate, in the coupled variant only \\
\midrule
\ref{var:sec:publicgood}
 & $X_t,P(t)$ & intracellular \yp{} count, and public-good concentration \\
 & $\lambda_0,\alpha$ & base birth rate and public-good advantage \\
 & $g,c,K^\ast$ & production, clearance, and the minimum founding cohort \cref{var:eq:kstar} \\
\midrule
\ref{var:sec:twophase}
 & $\lambda_1,\mu_1,\delta_1$ & phase-1 rates, of the intracellular
   birth--death--catastrophe process \\
 & $\lambda_2,\mu_2$ & phase-2 rates, of the extracellular birth--death
   process \\
 & $t_c$ & time of the switch between phases \\
 & $Z,p_0$ & offspring count of the host-to-host process, and the probability that the whole within-host process dies out \\
\bottomrule
\end{tabular}
\caption{Notation, organised by model. The rates $\lambda$, $\mu$ and
$\delta$, the roots $a$ and $b$ and the burst size $\Kb$ carry their
thesis-wide meanings everywhere in this chapter and are never redefined. Where
one letter serves two models --- $\alpha$ as an exponent in
\cref{var:sec:pimentel,var:sec:pimentelplus} and as a coupling coefficient in
\cref{var:sec:publicgood}, $V$ as potential in \cref{var:sec:rrrr} and as
vitality in \cref{var:sec:plague}, $c$ as a consumption rate in both, $N$ as
an abundance in \cref{var:sec:speciation} and as an integer capacity in
\cref{var:sec:pimentel} --- the two uses are separated by this table and never
occur in the same section.}
\label{var:tab:notation}
\end{table}

The ecological models occupy \cref{var:sec:speciation,var:sec:pimentel,%
var:sec:pimentelplus,var:sec:rrrr,var:sec:plague}. The \yp{} models close the
chapter: \cref{var:sec:publicgood} asks for the founding cohort that would
make an accumulating public good supercritical, and \cref{var:sec:twophase}
composes the two phases of plague infection and obtains $\ex{Z}$ in closed
form. \Cref{var:sec:discussion} gathers limitations and the open questions
that remain; \cref{var:app:models} summarises all seven models on one page,
\cref{var:app:figures} records how the figures were produced, and
\cref{var:app:pgf} gives the generating-function proof of
\cref{var:prop:yule-geom}.

\FloatBarrier
